All tables use \verb|booktabs| rules.
Never use vertical rules, and never use \verb|\hline| where \verb|\toprule|, \verb|\midrule|, or \verb|\bottomrule| will do.
The caption sits above the table, unlike a figure caption, which sits below.

\subsection{Simple table}
The plain case, with a column specification of \verb|l| for text and \verb|r| for numbers.
See \cref{tab:simple}.

\begin{table}[t]
  \centering
  \caption{A simple table with \texttt{booktabs} rules.}
  \label{tab:simple}
  \begin{tabular}{lrr}
    \toprule
    Technical process & Papers & Share \\
    \midrule
    System realisation      & 8281 & 96.02\% \\
    Verification            &  216 &  2.50\% \\
    Mission analysis        &  121 &  1.40\% \\
    Requirements definition &    5 &  0.06\% \\
    \bottomrule
  \end{tabular}
\end{table}

\subsection{Grouped headers and merged cells}
\verb|\multicolumn| groups columns, and \verb|\cmidrule(lr)| underlines only that group.
\verb|\multirow| merges cells downward.
The \verb|(lr)| trimming keeps adjacent rules from touching.
See \cref{tab:grouped}.

\begin{table}[t]
  \centering
  \caption{Grouped column headers with \texttt{\textbackslash multicolumn} and merged rows with \texttt{\textbackslash multirow}.}
  \label{tab:grouped}
  \begin{tabular}{llrr}
    \toprule
    \multirow{2}{*}{Phase} & \multirow{2}{*}{Process}
      & \multicolumn{2}{c}{Corpus} \\
    \cmidrule(lr){3-4}
      & & Count & Share \\
    \midrule
    \multirow{2}{*}{Concept}
      & Mission analysis   & 121 & 1.40\% \\
      & Stakeholder needs  &   0 & 0.00\% \\
    \midrule
    \multirow{2}{*}{Definition}
      & Requirements       &   5 & 0.06\% \\
      & Architecture       & 402 & 4.66\% \\
    \bottomrule
  \end{tabular}
\end{table}

\subsection{Table with a wrapping text column}
\verb|tabularx| stretches one column to fill a target width.
The class defines \verb|L|, \verb|C|, and \verb|R| as left, centre, and right ragged versions of the stretching \verb|X| column, which avoids the stretched justification that plain \verb|X| produces.
Set the target width to \verb|\columnwidth| in the two-column model.
See \cref{tab:wrapping}.

\begin{table}[t]
  \centering
  \caption{A wrapping description column, using \texttt{tabularx} and the \texttt{L} column type.}
  \label{tab:wrapping}
  \begin{tabularx}{\columnwidth}{lL}
    \toprule
    Process & Description \\
    \midrule
    Mission analysis
      & Establishes the problem space and the operational need that the system is meant to address. \\
    Requirements
      & Transforms stakeholder needs into a verifiable specification of what the system shall do. \\
    Verification
      & Confirms that the realised system meets the specification it was built against. \\
    \bottomrule
  \end{tabularx}
\end{table}

\subsection{Numeric alignment}
\verb|siunitx| aligns numbers on the decimal point through the \verb|S| column type.
Header text in an \verb|S| column must be wrapped in braces, or the package will try to parse it as a number.
Use \verb|\num| and \verb|\SI| in the body text too, so units and digit grouping stay consistent throughout the paper.
See \cref{tab:numeric}.

\begin{table}[t]
  \centering
  \caption{Decimal-aligned figures, using the \texttt{siunitx} \texttt{S} column.}
  \label{tab:numeric}
  \begin{tabular}{l S[table-format=4.2] S[table-format=2.3]}
    \toprule
    Protocol & {Rate (\si{\kilo\bit\per\second})} & {Loss (\si{\decibel})} \\
    \midrule
    BB84    & 1204.50 &  6.200 \\
    Decoy   &  980.75 & 11.480 \\
    MDI     &   12.05 & 24.000 \\
    TF      &    3.40 & 46.125 \\
    \bottomrule
  \end{tabular}
\end{table}

\subsection{Table with notes}
\verb|threeparttable| keeps footnotes attached to the table rather than to the page.
The note markers are independent of the document footnote counter, so they do not disturb the numbering of ordinary footnotes.
See \cref{tab:notes}.

\begin{table}[t]
  \centering
  \begin{threeparttable}
    \caption{A table carrying its own notes.}
    \label{tab:notes}
    \begin{tabular}{lrr}
      \toprule
      Source & Records\tnote{a} & Retained \\
      \midrule
      IEEE Xplore & 3412 & 3180 \\
      Scopus      & 6104 & 5444\tnote{b} \\
      \bottomrule
    \end{tabular}
    \begin{tablenotes}[flushleft]\footnotesize
      \item[a] Before de-duplication across the two databases.
      \item[b] Excludes conference abstracts shorter than two pages.
    \end{tablenotes}
  \end{threeparttable}
\end{table}

\subsection{Full-width table}
\verb|table*| spans both columns, and carries the same top-of-page-only restriction as \verb|figure*|.
See \cref{tab:wide}.

\begin{table*}[t]
  \centering
  \caption{A full-width table, declared with \texttt{table*}.}
  \label{tab:wide}
  \begin{tabularx}{\textwidth}{llLr}
    \toprule
    Process & Standard clause & Intent & Papers \\
    \midrule
    Mission analysis   & 6.4.1 & Characterise the problem space and the operational need. & 121 \\
    Stakeholder needs  & 6.4.2 & Elicit and agree the needs of those who will use the system. & 0 \\
    Requirements       & 6.4.3 & State verifiable system requirements derived from those needs. & 5 \\
    Architecture       & 6.4.4 & Define the structural options that could satisfy the requirements. & 402 \\
    \bottomrule
  \end{tabularx}
\end{table*}

\subsection{Oversized table}
When a table is a little too wide, shrink it with \verb|adjustbox| rather than reducing the font size by hand.
If it is far too wide, the table is usually trying to do too much, and belongs in an appendix or split in two.

\begin{table}[t]
  \centering
  \caption{A wide table clamped to the column with \texttt{adjustbox}.}
  \label{tab:clamped}
  \begin{adjustbox}{max width=\columnwidth}
    \begin{tabular}{lrrrrr}
      \toprule
      Process & 2021 & 2022 & 2023 & 2024 & 2025 \\
      \midrule
      Implementation & 412 & 455 & 501 & 588 & 640 \\
      Integration    &  88 & 104 & 161 & 224 & 302 \\
      Verification   &  21 &  24 &  30 &  36 &  44 \\
      \bottomrule
    \end{tabular}
  \end{adjustbox}
\end{table}
